\documentclass[a4paper]{article}
\usepackage[pdftex]{graphicx}
\usepackage[utf8]{inputenc}
\usepackage{enumerate}
\usepackage{siunitx}
\sisetup{locale=DE} 
\usepackage{href-ul}
\hypersetup{
	colorlinks=true,
	linkcolor=blue,
	urlcolor=blue}
\usepackage{icomma}
\usepackage{amssymb}
\usepackage{geometry}
\geometry{a4paper, top=15mm, left=15mm, right=15mm, bottom=15mm,
	headsep=10mm, footskip=12mm}

\begin{document}
{\bf Gravitation, Himmelsmechanik}
\begin{enumerate}[1.]

% Aufgabe 2
\begin{minipage}{0.7\textwidth}
{\bf \item Der Marsmond Phobos besitzt eine Masse von 
$\SI{1,1e16}{\kilogram}$ und einen Radius von etwa $\SI{11}{\kilo\meter}$\\}
$[\textnormal{Gravitationskonstante}~\gamma=\SI{6,67e-11}{\meter^3 \per {\kilogram\second^2}}]$
\end{minipage}
\hfill
\begin{minipage}{0.25\textwidth}
\includegraphics[width=2.5 cm]{phobos95}
\end{minipage}

\begin{enumerate}[a)]
\item Zeige durch Rechnung, dass auf der Oberfläche von Phobos eine Fallbeschleunigung von ca. $g_{Ph}=\SI{6,1e-3}{\meter\over {\second^2}}$ herrscht.
\item Auf  Phobos kann man mit sich selber Baseball spielen. Hierzu wirft man einen Ball horizontal ab, so dass dieser Phobos auf einer oberflächennahen Kreisbahn umrundet und wieder am Abwurfpunkt ankommt. Mit welcher Geschwindigkeit muss dieser Ball abgeworfen werden?
\item Nach dem Abwurf wartet man, bis der Ball am gegenüberliegenden Hori\-zont wieder auftaucht. Wie viel Zeit vergeht vom Abwurf bis zur Vollendung der Umrundung, also bis zum Abschlag?
\end{enumerate}


{\bf \item Auf der Erde ist der Ortsfaktor $g$ und der Radius} abhängig von der geographischen Breite. Der Erdradius ist am Pol $\SI{6357}{\kilo\meter}$, am Äquator $\SI{6378}{\kilo\meter}$, die Masse der Erde ist $\SI{5,97e24}{\kilogram}$

\begin{enumerate}[a)]
	\item Welche zwei Effekte sind für den Unterschied von $g_{Pol}-g_{Aeq}$ verantwortlich?
	\item Berechne den absoluten und den relativen Unterschied von $g_{Pol}$ und $g_{Aeq}$
\end{enumerate}

{\bf \item Schwarze Löcher haben ein so starkes Gravitationsfeld}, dass nicht einmal mehr Licht schnell genug ist, um zu entweichen, wenn es sich näher als der Schwarzschild-Radius befindet. Berechne diesen Radius für ein Schwarzes Loch mit 10-facher Sonnenmasse $M_{Sonne}= \SI{1,98e30}{\kilogram}$

{\bf \item Zwischen Mond und Erde gibt es einen Punkt,} in dem sich die Gravitationsfeldstärken gegenseitig aufheben. Berechne den Abstand $r$ dieses Punktes zum Mondmittelpunkt. Der Abstand Mond-Erde sei $d=\SI{384400}{\kilo\meter}$. Masse Mond: $\SI{7,35e22}{\kilogram}$ 

\includegraphics[width=18 cm]{moe413}
	{\bf \item Raumfahrer Spiff ist mit seiner Landefähre in einen Orbit um einen Exoplaneten geschwenkt}.
\noindent
\begin{minipage}{0.4\textwidth}
	Er misst folgende Daten: 
	\begin{enumerate}[a)]
		\item Wie groß ist der Radius des Planeten?
		\item Wie groß ist die Masse des Planeten?
		\item Welche Gravitationsfeldstärke wird Spiff bei einer Landung erwarten?
	\end{enumerate}
	
\end{minipage}
\hspace{1 cm}
\begin{minipage}{0.4\textwidth}
 \renewcommand{\arraystretch}{1.5}
 \begin{tabular}{|c|c|}
 	\hline
 	Höhe über Grund & $h=\SI{800}{\kilo\meter}$ \\
 	\hline
 	Geschwindigkeit über Grund & $v=\SI{3,7}{\kilo\meter \over \second}$\\
 	\hline
 	Umlaufdauer & $T=\SI{130}{\minute}$\\
 	\hline
 \end{tabular}
\end{minipage}

 [Die Eigenrotation des Exoplaneten ist zu vernachlässigen!]	





\end{enumerate} 

\fbox{
	\begin{minipage}{0.5\textwidth}
		Zur Lösung bitte \href{https://www.okuyakl.de/phys/p10kreiL413/ll413.pdf}{hier klicken} oder den QR-Code scannen.\\
	Weitere Arbeitsblätter gibt es unter 
	
	\href{https://www.okuyakl.de}{www.okuyakl.de}
	\end{minipage}
	\hfill
	\begin{minipage}{0.4\textwidth}
		\includegraphics[width=1.5 cm]{../../viecher/zwe03}
		\includegraphics[width=3 cm]{qr413}
		\includegraphics[width=2 cm]{../../viecher/afanticon1}
		
	\end{minipage}}
	
\end{document}